\section{Перевірка виконання загального завдання}

У цьому розділі наведено детальний опис та аналіз практичної реалізації 
перших трьох пунктів програми досліджень, орієнтованих на освоєння базових 
інтерфейсів \textit{Java Collections Framework}.

\subsection{Робота зі списками об'єктів (\texttt{List<Student>})}

Згідно з пунктом 1 завдання, було спроектовано динамічний лінійний список 
на базі узагальненого класу \texttt{ArrayList<Student>}
. На відміну 
від статичних масивів, дана структура дозволяє гнучко додавати нові об'єкти 
користувацького класу без ручного контролю ємності пам'яті. 

В межах тестування було виконано такі операції:
\begin{itemize}
    \item \textbf{Динамічне розширення:} 
    За допомогою вбудованого методу \texttt{add()} до списку було послідовно 
    додано трьох студентів із фіксованими анкетними даними.
    \item \textbf{Ітераційний аналіз даних:} 
    За допомогою циклу \texttt{for-each} реалізовано прозорий перебір елементів 
    колекції. Застосування методу \texttt{size()} дозволило динамічно 
    розрахувати середній академічний бал групи, повністю нівелюючи ризики 
    помилок «жорсткого кодування» лімітів масивів.
\end{itemize}

\subsection{Робота з унікальними множинами (\texttt{Set<String>})}

Пункт 2 завдання присвячений дослідженню механізмів гарантії унікальності даних 
за допомогою інтерфейсу \texttt{Set} на прикладі реалізації \texttt{HashSet<String>}. 

Для демонстрації захисту від появи дублікатів (аналог обмеження \texttt{UNIQUE} 
у реляційних СУБД) до колекції було виконано чотири спроби додавання назв навчальних 
дисциплін, причому ідентифікатор \texttt{"Java"} надсилався до сховища двічі. 
Результати тестування повністю підтвердили низькорівневу логіку роботи хеш-таблиці: 
дублікат був автоматично відфільтрований за допомогою методів \texttt{hashCode()} 
та \texttt{equals()}, і у множині зберігся лише один унікальний екземпляр 
рядка.

\subsection{Реалізація індексної таблиці відображень (\texttt{Map<Integer, String>})}

Згідно з пунктом 3 завдання, було розгорнуто асоціативну карту 
\texttt{HashMap<Integer, String>}, яка виступає прямим архітектурним аналогом 
організації індексного пошуку в базах даних за \textbf{первинним ключем (Primary Key)}.

У створеній структурі як унікальний ключ (\textit{Key}) використовується 
обгортка примітиву \texttt{Integer} (код студента), а як значення 
(\textit{Value}) — об'єктний рядок \texttt{String} (ПІБ). Практичне 
тестування підтвердило константну швидкість доступу до даних за допомогою методу 
\texttt{get(102)}, а також продемонструвало ефективність повного перебору 
пар «ключ -- значення» через ітераційну множину записів \texttt{entrySet()}.
